\scanpage{30}
% archive sheet number: 29

La connaissance de l'ensemble $X$ et de la famille des corps $(k(x))_{x\in X}$ (corps résiduels \deleted{\unsure{des}} en les pts de $\mathrm{Spec}\,A$ ou $\mathrm{Spec}(\mathscr X)$ -- $k(x)$ représente le foncteur $k'$ (corps variable sur $k$) $\longmapsto$ ens des \deleted{points} éléments de $\mathscr X(k')$ ayant pour lieu $x$).

\textcircled{b} \uline{Topologie sur $X$} (les fermés sont les) $V(S)\subset X$ associés aux parties $S\inserted{=(f_\alpha)}$ de $A$, \deleted{\unsure{leur}} un point \inserted{géom. $u$} à valeurs dans $k'$ a son lieu dans $V(S)$ sss $f_\alpha(u)=0$ pour tout $\alpha$. On a \deleted{$V(S)=V(J_S)$ si} $V(S)=$ ensemble des idéaux premiers de $A$ qui contiennent $S$.

On a $V(S)=V(J_S)$, $J_S$ idéal engendré, et $\unsure{\overline{\mathrm m}|=}$ $V(S)=V(\widetilde J_S)$. Il \deleted{\unsure{On a $V($|Donc}} est bon de se borner aux idéaux, et $\overline{\mathrm m}$ aux idéaux ``radicaux''. \deleted{L'ensemble des} On voit que $V(\unsure{S|J})=V(\unsure{S'|J'})$ (resp. $V(S)\subset V(S')$) signifie, avec des notations précédentes, que $V'_S=V^{\unsure{\prime\prime|\prime}}_{S'}$ (restrictions aux corps) \deleted{du moins que} (resp $V'_S\deleted{\unsure{(k)}}\subset V'_{S'}$), donc $\widetilde J_S=\widetilde J_{S'}$ (resp. $\widetilde J_S\supset\widetilde J_{S'}$) : cela provient au fond que $\widetilde J_S$ est l'intersection des idéaux premiers qui contiennent $A$.

\[
\deleted{V(\sum J_i}\ V(\bigcup_{i\in I}S_i)=V(\sum J_i)=\bigcap_i V(S_i)=\bigcap V(J_i)
\qquad I\ \text{quelc.}
\]

(NB une \deleted{intersection} somme de deux \inserted{idéaux} \deleted{\unsure{rad}} radicaux n'est pas nécessairement un idéal radical)

\[
V(\bigcap_i J_i)=V(\prod J_i)=\bigcup_i V(J_i)\qquad I\ \text{fini}
\]

[NB $\prod_i J_i\subset\bigcap_i J_i$, il faut prouver que $\forall$ idéal premier $\mathfrak p$ qui contient l'\deleted{intersection}\inserted{produit} des $J_i$ contient un des $J_i$, or sinon on aurait pour $\forall\ i\in I$ un $f_i\in J_i,\notin\mathfrak p$, \deleted{on en passant à $A/\mathfrak p$ que si un produit des idéaux} donc $\prod f_i$ serait $\in\bigcap J_i$ et $\notin\mathfrak p$, absurde ] \deleted{\unsure{premiers est nul}, l'un d'eux est nul}

Donc on a une \uline{topologie}.

\margin{Si $x\in X$ correspond à $p_x\subset A$, Alors $\overline{\{x\}}$ est le plus petit fermé $V(J)$ qui contient $x$, i.e. tel que $J\subset p_x$, donc \deleted{\unsure{plus}} c'est $V(p_x)$.}

\uline{Espace \deleted{\unsure{top.}} quasi-compact} \inserted{(mais non séparé)}

\margin{$x\in\overline{\{y\}}$ $\Updownarrow$ $\overline{\{x\}}\subset\overline{\{y\}}\Longleftrightarrow p_x\supset p_y$}

$[\,X_f=X-V(f)$ est ouvert, \inserted{(le lieu de $u\in\mathscr X(k)$ est dans $X_f$ sss $f(u)\ne0$)}, et les \deleted{\unsure{$D(f)$}} $X_f$ forment une base de la topologie de $X$, stable par intersections finies ($X_{1}=X$, $X_{fg}=X_f\cap X_g$).

% slanted
\margin{plus bas}
